\section{Vorticity-based interpretation of dimensional effects}
\label{sec:uej-vorticity-dynamics}

The preceding comparisons show that the axisymmetric cylindrical and
three-dimensional Cartesian formulations give similar mean wave structures
over approximately the first two shock cells, but different downstream
behaviour. The mean-density modulation in the axisymmetric calculation weakens
rapidly after the second cell, while its centreline density subsequently
approaches the ambient value too slowly. The Cartesian calculation retains
distinct compression and expansion regions farther downstream and eventually
approaches the ambient state more closely. This section examines whether these
differences are consistent with the vorticity dynamics permitted by the two
formulations. The diagnostics are based on selected instantaneous fields and
are used to develop a physical interpretation rather than a closed or
statistically converged vorticity budget.

\paragraph{Compressible vorticity transport.}
In the absence of non-conservative body forces, the vorticity
\(\boldsymbol{\omega}=\nabla\times\boldsymbol{u}\) satisfies
\begin{equation}
    \frac{D\boldsymbol{\omega}}{Dt}
    =
    \underbrace{(\boldsymbol{\omega}\cdot\nabla)\boldsymbol{u}}
        _{\boldsymbol{T}_s}
    -\underbrace{\boldsymbol{\omega}\,\nabla\cdot\boldsymbol{u}}
        _{\boldsymbol{T}_d}
    +\underbrace{\frac{\nabla\rho\times\nabla p}{\rho^2}}
        _{\boldsymbol{T}_b}
    +\underbrace{\nabla\times\!\left[
        \frac{1}{\rho}\nabla\cdot\boldsymbol{\tau}\right]}
        _{\boldsymbol{T}_v}.
    \label{eq:uej-vorticity-transport}
\end{equation}
The terms on the right represent stretching and tilting, dilatation,
baroclinic torque, and viscous transport, respectively. Comparable
decompositions have been used to examine vorticity transport in supersonic
jets \citep{Li2016,BhideCuppoletti2024}. The stretching term changes the
magnitude and direction of existing vorticity. The dilatation term amplifies
vorticity in compression and weakens it in expansion. The baroclinic term is
the inviscid volumetric source that can generate vorticity from an initially
irrotational state when the pressure and density gradients are not aligned.
It vanishes across an ideal locally one-dimensional shock for which these
gradients are parallel, but it can be important where a shock intersects a
shear layer or another non-uniform region. The viscous term can also generate,
redistribute, and dissipate vorticity. Its constant-property reduction to
\(\nu\nabla^2\boldsymbol{\omega}\) is not applicable here because density and
the temperature-dependent transport properties vary through the flow.

For a captured shock, the physical shock thickness is smaller than the grid
spacing. Derivatives evaluated inside the numerical shock therefore depend on
the discrete shock structure. The following comparison emphasises spatial
distributions and radially integrated diagnostics rather than isolated local
maxima within a shock.

\paragraph{Axisymmetric restriction.}
For an axisymmetric flow without swirl,
\(\partial/\partial\theta=0\), \(u_\theta=0\), and the vorticity contains only
the azimuthal component,
\[
    \boldsymbol{\omega}=\omega_\theta\boldsymbol{e}_\theta,
    \qquad
    \omega_\theta
    =\frac{\partial u_r}{\partial x}
    -\frac{\partial u_x}{\partial r}.
\]
Its transport equation is
\begin{equation}
    \frac{D\omega_\theta}{Dt}
    =
    \frac{u_r}{r}\omega_\theta
    -\omega_\theta\left(
        \frac{\partial u_x}{\partial x}
        +\frac{\partial u_r}{\partial r}
        +\frac{u_r}{r}\right)
    +\frac{1}{\rho^2}\left(
        \frac{\partial\rho}{\partial x}\frac{\partial p}{\partial r}
        -\frac{\partial\rho}{\partial r}\frac{\partial p}{\partial x}
      \right)
    +T_{v,\theta}.
    \label{eq:uej-axisymmetric-vorticity}
\end{equation}
The first term in Equation~\eqref{eq:uej-axisymmetric-vorticity} is the
geometric stretching of a vortex ring as its radius changes. Axisymmetry does
not therefore remove vortex stretching altogether. It does, however, exclude
azimuthal deformation and the tilting of circumferential vorticity into axial
and radial components. Axisymmetric vortex rings can roll up, pair, and move
radially, but they cannot access the full component redistribution available
to the Cartesian flow.

\paragraph{Enstrophy production and comparison protocol.}
Projecting Equation~\eqref{eq:uej-vorticity-transport} onto
\(\boldsymbol{\omega}\) gives
\begin{equation}
    \frac{D\mathcal{E}_\omega}{Dt}
    =P_s+P_d+P_b+P_v,
    \qquad
    \mathcal{E}_\omega=\frac{1}{2}|\boldsymbol{\omega}|^2,
    \label{eq:uej-enstrophy-transport}
\end{equation}
where
\begin{equation}
\begin{aligned}
    P_s&=\boldsymbol{\omega}\cdot
        (\boldsymbol{\omega}\cdot\nabla)\boldsymbol{u},\\
    P_d&=-|\boldsymbol{\omega}|^2\nabla\cdot\boldsymbol{u},\\
    P_b&=\frac{\boldsymbol{\omega}\cdot
        (\nabla\rho\times\nabla p)}{\rho^2},\\
    P_v&=\boldsymbol{\omega}\cdot\boldsymbol{T}_v.
\end{aligned}
    \label{eq:uej-enstrophy-terms}
\end{equation}
The signed inviscid contributions are normalised as
\(P_k^*=P_kD_e^3/U_j^3\), with \(k\in\{s,d,b\}\). In the axisymmetric
formulation,
\(P_s=(u_r/r)\omega_\theta^2\),
\(P_d=-\omega_\theta^2\nabla\cdot\boldsymbol{u}\), and
\(P_b=\omega_\theta(\nabla\rho\times\nabla p)_\theta/\rho^2\).
The viscous contribution requires second derivatives and is more sensitive to
AMR interfaces; it is not included in the present diagnostic comparison.

The axisymmetric Level~4 and Cartesian Level~5 calculations have the same
nominal finest spacing, \(D_e/256\), but their refinement footprints are not
identical. As shown in Table~\ref{tab:uej-refinement-criteria}, the finest
Cartesian region is concentrated near the nozzle lip, whereas the
axisymmetric Level~4 region covers a larger part of the near-field core. The
comparison is therefore not locally resolution matched throughout the plotted
window. The fields are evaluated at \(t=8.0\)~ms for the axisymmetric case and
\(t=10.16\)~ms for the Cartesian case, after the initial transient in each
calculation. Derivatives are evaluated separately on the native AMR levels and
assembled into a finest-owned composite without differentiating across a
coarse--fine interface. The Cartesian diagnostics are evaluated on the
\(z=0\) plane. The two snapshots do not represent a phase-matched pair, so the
comparison is treated as an exploratory instantaneous assessment.

\paragraph{Vorticity components.}
Figure~\ref{fig:uej-dim-vortcomp} compares numerical schlieren, the signed
azimuthal vorticity, and the magnitude of the axial and radial components. The
axisymmetric snapshot retains large coherent toroidal structures, whereas the
Cartesian shear layer contains smaller and less circumferentially coherent
structures downstream. Axial and radial vorticity, which vanish identically in
the axisymmetric formulation, become visible in the Cartesian shear layer from
approximately \(x/D_e=0.6\). Over the plotted Cartesian field, the spatial
99th percentiles of
\(\sqrt{\omega_x^2+\omega_r^2}\) and \(|\omega_\theta|\) are approximately
\(19U_j/D_e\) and \(17U_j/D_e\), respectively. The substantial
non-azimuthal content is consistent with redistribution among vorticity
components, but does not by itself distinguish tilting from baroclinic
generation.

\begin{figure}[!htbp]
    \centering
    \includegraphics[width=\textwidth]{paperfig_dim_vortcomp.pdf}
    \caption{Instantaneous vorticity structure in the axisymmetric
    cylindrical Level~4 calculation (upper row) and on the \(z=0\) plane of
    the three-dimensional Cartesian Level~5 calculation (lower row):
    numerical schlieren (left), signed azimuthal vorticity (centre), and the
    magnitude of the axial and radial vorticity components (right). The
    axisymmetric field is reflected about the axis. The two hierarchies have
    the same nominal finest spacing, \(D_e/256\), but different spatial
    refinement footprints.}
    \label{fig:uej-dim-vortcomp}
\end{figure}

\paragraph{Production terms.}
Figure~\ref{fig:uej-dim-production} shows the three signed inviscid enstrophy
terms on the same symmetric logarithmic scale. The largest difference is in
\(P_s^*\). Its spatial 99th-percentile magnitude is approximately
\(1.42\times10^3\) in the Cartesian snapshot and 81 in the axisymmetric
snapshot. The Cartesian stretching field becomes distributed across the
developing shear layer, while the axisymmetric geometric contribution changes
sign between neighbouring radial structures. By comparison, the
99th-percentile magnitudes of \(P_d^*\) are approximately 308 and 280 in the
Cartesian and axisymmetric snapshots, and those of \(P_b^*\) are approximately
79 and 91. Dilatation and baroclinic activity therefore have comparable
instantaneous orders of magnitude in the two selected fields, while the
stretching contribution differs much more strongly. For highly underexpanded
jets, \citet{Li2016} also identified stretching as an important mechanism in
the distortion of vortex rings. The present result applies only to the
moderately underexpanded condition \(\eta_0=3.2735\)
\((\eta_e\simeq1.73)\) and to the selected snapshots.

\begin{figure}[!htbp]
    \centering
    \includegraphics[width=\textwidth]{paperfig_dim_production.pdf}
    \caption[Instantaneous inviscid enstrophy contributions in the
    axisymmetric and Cartesian calculations.]{Instantaneous normalised
    enstrophy contributions in the
    axisymmetric cylindrical Level~4 calculation (upper row) and on the
    \(z=0\) plane of the three-dimensional Cartesian Level~5 calculation
    (lower row): stretching \(P_s^*\) (left), dilatation \(P_d^*\) (centre),
    and baroclinic contribution \(P_b^*\) (right). All panels use the same
    symmetric logarithmic colour scale.}
    \label{fig:uej-dim-production}
\end{figure}

\paragraph{Plane-based radial estimates.}
For each contribution, the dimensionless net and activity measures are
defined as
\begin{equation}
\begin{aligned}
    \Pi_k^{\mathrm{net}}(x)
        &=\frac{D_e}{U_j^3}
          \int_0^{r_{\max}}P_k^{(p)}(x,r)\,2\pi r\,\mathrm{d}r,\\
    \Pi_k^{\mathrm{act}}(x)
        &=\frac{D_e}{U_j^3}
          \int_0^{r_{\max}}|P_k^{(p)}(x,r)|\,2\pi r\,\mathrm{d}r.
\end{aligned}
    \label{eq:uej-plane-production-measures}
\end{equation}
For the axisymmetric calculation, \(P_k^{(p)}\) is the computed radial field
and the ring weighting is exact. For the Cartesian calculation, it is the
average of the positive and negative \(y\) traces on the \(z=0\) plane. The
Cartesian result is therefore an azimuthally reconstructed plane-based
estimate, not a full three-dimensional cross-sectional integral. The plotted
curves are smoothed over \(0.1D_e\) in the axial direction for display.

Figure~\ref{fig:uej-dim-integrals}(a) shows a predominantly positive
Cartesian stretching estimate between approximately \(x/D_e=1\) and 4,
whereas the corresponding axisymmetric net value remains much smaller because
positive and negative contributions cancel across the sampled radial
cross-section. The activity measure in panel (b) confirms that the small net
axisymmetric value does not imply that the local terms vanish. At
\(x/D_e=2\), the Cartesian stretching activity in this snapshot is
approximately 50 times the axisymmetric value. This factor describes the
selected instantaneous plane estimates and is not interpreted as a
time-averaged production ratio.

The redistribution of vorticity components is summarised by
\begin{equation}
    F_{\mathrm{3D}}(x)
    =\frac{
       \displaystyle\int_0^{r_{\max}}
       \left(\omega_x^2+\omega_r^2\right)^{(p)}2\pi r\,\mathrm{d}r}
      {
       \displaystyle\int_0^{r_{\max}}
       \left(|\boldsymbol{\omega}|^2\right)^{(p)}2\pi r\,\mathrm{d}r}.
    \label{eq:uej-three-dimensional-vorticity-fraction}
\end{equation}
The quantity is identically zero in the axisymmetric formulation. In the
Cartesian plane estimate it exceeds one half near \(x/D_e=1.1\), reaches
approximately 0.77 near \(x/D_e=2\), and subsequently fluctuates around values
close to two thirds. Two thirds is the component-equipartition reference for
locally isotropic vorticity. Proximity to this reference shows that
non-azimuthal components have become substantial, but does not by itself
establish local isotropy.

\begin{figure}[!htbp]
    \centering
    \includegraphics[width=\textwidth]{paperfig_dim_integrals.pdf}
    \caption{Plane-based radial estimates of the instantaneous enstrophy
    terms: (a) signed net contribution, (b) activity obtained by integrating
    the absolute value, and (c) the fraction of axial and radial vorticity in
    the Cartesian calculation. Solid lines denote the Cartesian Level~5
    snapshot and dashed lines the axisymmetric Level~4 snapshot. The Cartesian
    curves use the average of the two radial traces on the \(z=0\) plane with
    an axisymmetric ring weight and are not full cross-sectional integrals.
    Curves in panels (a) and (b) are smoothed over \(0.1D_e\) for display. The
    dotted line in panel (c) marks the component-equipartition reference of
    two thirds.}
    \label{fig:uej-dim-integrals}
\end{figure}

Taken together, the three diagnostics are consistent with the interpretation
that the additional vorticity-component pathways available in the Cartesian
formulation promote the loss of circumferential coherence and alter the
subsequent mixing of the jet. This interpretation is also consistent with the
different downstream mean-density behaviour reported in
Sections~\ref{subsec:uej-panda-centreline}
and~\ref{subsec:uej-panda-radial}. The comparison does not, however, isolate
dimensionality as the sole cause: the two diagnostics use different
instantaneous states and non-identical AMR footprints, and the viscous
contribution is not included. It therefore supports a vorticity-based
interpretation of the dimensional difference rather than a closed causal
proof. The similar mean wave structure over approximately the first two shock
cells remains the appropriate range in which the axisymmetric formulation is
used as a lower-cost reference.

\FloatBarrier
